% strat_t4_subperiod — subperiod performance of the GCF bond-timing strategy.
% Numbers transcribed from the output of R/strategy_ext.R.
\small
\begin{tabular}{l c c c}
  \toprule
  Strategy & Mean (\%) & Vol (\%) & Sharpe \\
  \midrule
  \panel{4}{Panel A: 2005--2014} \\
  $\GCF$-timed           & $6.47$ & $6.98$ & $0.93$ \\
  Recursive-mean timing  & $5.47$ & $5.37$ & $1.02$ \\
  Buy-and-hold           & $5.33$ & $4.90$ & $1.09$ \\
  \midrule
  \panel{4}{Panel B: 2015--2024} \\
  $\GCF$-timed           & $-0.34$ & $7.28$ & $-0.05$ \\
  Recursive-mean timing  & $-0.61$ & $7.12$ & $-0.09$ \\
  Buy-and-hold           & $-0.91$ & $7.34$ & $-0.12$ \\
  \midrule
  \panel{4}{Panel C: 2021--2022 (hiking cycle)} \\
  $\GCF$-timed           & $-9.83$  & $6.29$ & $-1.56$ \\
  Recursive-mean timing  & $-9.72$  & $6.04$ & $-1.61$ \\
  Buy-and-hold           & $-10.49$ & $6.02$ & $-1.74$ \\
  \bottomrule
\end{tabular}
\tabnotes{This table splits the out-of-sample window of \Cref{tab:strat-perf}
  into subperiods. Panels A and B are the two halves of the evaluation window;
  Panel C isolates the 2021--2022 global hiking cycle, the worst bond-market
  episode in the sample. Strategies, exposure rule, and equal-average-exposure
  reporting are identical to \Cref{tab:strat-perf}; returns are local-currency
  (currency-hedged) excess returns based on overlapping twelve-month holding
  periods, so subperiod Sharpe ratios are descriptive rather than tested.}
